%%%%%%%%%%%%%%%%%%%%%%%%%%%%%%%%%%%%%%%%%%%%%%%%%%%%%%%%%%%%%%%%%%%%%%%%%%%%%%%
%                         CAPÍTULO 2
%                  Redes Inalámbricas y Comunicaciones IoT
%%%%%%%%%%%%%%%%%%%%%%%%%%%%%%%%%%%%%%%%%%%%%%%%%%%%%%%%%%%%%%%%%%%%%%%%%%%%%%%

\chapter{Redes Inalámbricas y Comunicaciones IoT}

En los capítulos anteriores hemos estudiado cómo transmitir datos mediante cables:
señales eléctricas que viajan por hilos de cobre siguiendo reglas estrictas como
UART o RS-232. Sin embargo, en muchos entornos industriales y de automatización
resulta inviable o poco práctico tender cables entre cada dispositivo. Imagina un
almacén robotizado donde cientos de sensores de movimiento deben comunicarse con
una central, o una fábrica de alimentos donde los cables dificultan la higiene.

En este capítulo abordamos las \textbf{comunicaciones inalámbricas}, que utilizan
ondas de radio en lugar de conductores físicos para transmitir información. Nos
centraremos en las dos tecnologías más utilizadas en sistemas embebidos:
\textbf{Bluetooth Low Energy} (BLE) y \textbf{Wi-Fi}, y veremos cómo aplicarlas
en la práctica para construir una red de sensores IoT (Internet de las Cosas).


\section{Comunicaciones inalámbricas: ¿por qué eliminar los cables?}

La tendencia en la industria 4.0 es hacia la \textbf{desconexión física}: dispositivos
que se comunican sin cables, que se autoconfiguran y que pueden desplazarse.
Algunas ventajas de las comunicaciones inalámbricas son:

\begin{itemize}
    \item \textbf{Flexibilidad de instalación}: no es necesario diseñar canalizaciones
    ni prever la longitud exacta de los cables.
    \item \textbf{Movilidad}: robots, carretillas elevadoras o sensores portátiles pueden
    cambiar de ubicación sin restringirse por el cableado.
    \item \textbf{Reducción de costes de mantenimiento}: los cables físicos se degradan
    con el tiempo, especialmente en entornos con vibraciones, humedad o
    temperaturas extremas.
    \item \textbf{Escalabilidad}: es más sencillo añadir nuevos dispositivos a una red
    inalámbrica existente que tender nuevos cables.
\end{itemize}

Sin embargo, también existen inconvenientes:
\begin{itemize}
    \item \textbf{Interferencias}: otras redes, microondas o motores eléctricos pueden
    degradar la señal.
    \item \textbf{Seguridad}: las ondas de radio pueden ser interceptadas si no se
    cifran adecuadamente.
    \item \textbf{Consumo energético}: transmitir por radio consume más energía que
    enviar señales por un hilo.
    \item \textbf{Alcance limitado}: la potencia de transmisión está regulada y, en
    interiores, las paredes atenúan la señal.
\end{itemize}


\section{Bluetooth vs. Wi-Fi: dos filosofías, una misma banda}

Tanto Bluetooth como Wi-Fi operan en la banda de radio de \textbf{2.4 GHz} (aunque
Wi-Fi también puede usar 5 GHz y 6 GHz en versiones modernas). A pesar de
compartir el mismo espectro, están diseñados para propósitos muy diferentes.

\begin{definicion}{Bluetooth y Wi-Fi}
\textbf{Bluetooth} es un estándar de comunicación inalámbrica de corto alcance y
bajo consumo, pensado para crear redes de área personal (PAN) entre dispositivos
próximos.

\textbf{Wi-Fi} es un estándar de comunicación inalámbrica de mayor alcance y
velocidad, diseñado para crear redes de área local (LAN) que permiten a
múltiples dispositivos acceder a internet o compartir recursos de forma simultánea.
\end{definicion}

\begin{table}[htbp]
\centering
\small
\caption{Comparación entre Bluetooth y Wi-Fi}
\begin{tabularx}{\textwidth}{@{}lXX@{}}
\toprule
\textbf{Característica} & \textbf{Bluetooth (BLE)} & \textbf{Wi-Fi} \\
\midrule
Alcance típico & 10--100 m (según clase) & 30--100 m (interior) \\
Velocidad de datos & Hasta 2 Mb/s (BLE 5.0) & Hasta varios Gb/s (Wi-Fi 6) \\
Consumo de energía & Muy bajo (meses con pila) & Alto (necesita alimentación) \\
Topología típica & Punto a punto (1 maestro, 7 esclavos) & Estrella (AP + múltiples clientes) \\
Uso principal & Periféricos, sensores, wearables & Acceso a internet, streaming, redes corporativas \\
Complejidad de protocolo & Simple (GATT) & Compleja (TCP/IP, HTTP, etc.) \\
Interferencias en 2.4 GHz & Alta (comparte banda con Wi-Fi) & Alta (comparte banda con Bluetooth) \\
\bottomrule
\end{tabularx}
\end{table}

\textbf{¿Cuándo usar cada uno?}
\begin{itemize}
    \item Usa \textbf{Bluetooth} cuando necesites conectar un sensor o un periférico
    a un dispositivo central, con bajo consumo y sin necesidad de acceso a internet.
    Ejemplos: auriculares inalámbricos, pulsómetros, sensores de temperatura portátiles.
    \item Usa \textbf{Wi-Fi} cuando el dispositivo necesite enviar datos a un servidor
    en la nube, recibir comandos remotos o formar parte de una red corporativa.
    Ejemplos: cámaras de videovigilancia, termostatos inteligentes, paneles de control.
\end{itemize}


\section{Bluetooth Low Energy (BLE)}

\subsection{¿Qué es BLE?}

Bluetooth Low Energy (BLE), también conocido como Bluetooth Smart, es una versión
del estándar Bluetooth optimizada para un consumo de energía mínimo. No es
compatible a nivel de protocolo con el Bluetooth clásico (el de los auriculares), pero
comparte la misma banda de radio.

\begin{definicion}{Bluetooth Low Energy (BLE)}
BLE es un protocolo de comunicación inalámbrica de corto alcance que permite la
transmisión de pequeños paquetes de datos con un consumo de energía extremadamente
bajo. Es ideal para dispositivos alimentados por baterías que deben operar durante
meses o años sin recargarse.
\end{definicion}

En BLE, los dispositivos desempeñan dos roles principales:
\begin{itemize}
    \item \textbf{Periférico (Peripheral)}: dispositivo que ofrece datos. Suele ser el sensor
    o el actuador.
    \item \textbf{Central (Central)}: dispositivo que solicita y recibe datos. Suele ser un
    teléfono, un PC o un gateway IoT.
\end{itemize}

\subsection{Arquitectura GATT: servicios y características}

BLE organiza la información mediante una estructura jerárquica conocida como
\textbf{GATT} (Generic Attribute Profile):

\begin{enumerate}
    \item \textbf{Perfil (Profile)}: conjunto de servicios predefinidos para un uso específico
    (por ejemplo, un perfil de pulsómetro).
    \item \textbf{Servicio (Service)}: agrupación lógica de características. Por ejemplo, un
    ``servicio de temperatura''.
    \item \textbf{Característica (Characteristic)}: el dato concreto. Puede ser una
    temperatura, una cadena de texto o un valor de configuración.
    \item \textbf{Descriptor (Descriptor)}: metadatos adicionales sobre una característica,
    como su formato o unidades.
\end{enumerate}

Cada servicio y característica se identifica mediante un \textbf{UUID} (Universally Unique
Identifier), un código de 128 bits que garantiza que no haya colisiones entre
aplicaciones diferentes.

\begin{definicion}{UUID en BLE}
Un \textbf{UUID} (Universally Unique Identifier) es un identificador numérico de 128 bits
que se utiliza para distinguir de forma única servicios, características y descriptores
dentro de una aplicación BLE. Los UUIDs pueden ser estandarizados (16 o 32 bits) o
customizados por el desarrollador (128 bits).
\end{definicion}

\subsection{Advertising y escaneo}

Un periférico BLE no puede ser conectado directamente. Primero debe anunciar su
presencia mediante un proceso llamado \textbf{advertising} (difusión): envía paquetes
periódicos de anuncio que contienen su nombre, el UUID de sus servicios y otros
datos.

Una central que desea conectarse realiza un \textbf{escaneo} (scanning): escucha los
paquetes de anuncio hasta encontrar un dispositivo con el UUID que busca. Una vez
encontrado, se establece la conexión.

\subsection{Notificaciones y lecturas}

Las características BLE pueden tener diferentes propiedades:
\begin{itemize}
    \item \textbf{READ}: la central puede leer el valor explícitamente.
    \item \textbf{WRITE}: la central puede escribir un valor.
    \item \textbf{NOTIFY}: el periférico envía el valor automáticamente cada vez que cambia,
    sin que la central tenga que pedirlo.
\end{itemize}

Las notificaciones son especialmente útiles en IoT: el sensor envía una nueva lectura
de temperatura cada vez que está disponible, sin que el receptor tenga que estar
constantemente preguntando.


\section{Wi-Fi en sistemas embebidos}

\subsection{Modos de operación}

Un dispositivo Wi-Fi puede funcionar en dos modos principales:

\begin{itemize}
    \item \textbf{STA (Station)}: el dispositivo se conecta a un punto de acceso (AP)
    existente, como un router. Es el modo más común para sensores IoT.
    \item \textbf{AP (Access Point)}: el dispositivo crea su propia red Wi-Fi a la que
    otros dispositivos pueden conectarse. Útil para configuración inicial.
\end{itemize}

\subsection{Protocolo HTTP para IoT}

En la práctica, cuando un dispositivo embebido con Wi-Fi necesita comunicarse, lo
hace a través del protocolo \textbf{HTTP} (HyperText Transfer Protocol), el mismo que
utiliza un navegador web para cargar páginas.

\begin{definicion}{HTTP en IoT}
HTTP es el protocolo de comunicación que permite a un cliente (por ejemplo, una
ESP32 o un navegador) solicitar recursos a un servidor. En el contexto IoT, un sensor
puede actuar como \textbf{servidor HTTP}, respondiendo a peticiones con datos en formato
texto, JSON o XML.
\end{definicion}

El flujo típico es:
\begin{enumerate}
    \item El dispositivo sensor se conecta a la red Wi-Fi (modo STA).
    \item Se le asigna una dirección IP por DHCP.
    \item Inicia un servidor HTTP en un puerto (normalmente el 80).
    \item Define rutas, como \texttt{/temperatura}, que devuelven el valor del sensor.
    \item Un cliente (otra ESP32, un PC o un smartphone) realiza una petición
    \texttt{GET} a esa IP y ruta, y recibe el dato.
\end{enumerate}

\subsection{Redes corporativas y seguridad}

En entornos universitarios o industriales, las redes Wi-Fi suelen estar protegidas
mediante autenticación (WPA2-PSK, WPA2-Enterprise) o listas de control de acceso
(MAC filtering). Esto puede complicar la conexión de dispositivos IoT, ya que no todos
los sistemas embebidos soportan los métodos de autenticación más avanzados.

En la práctica, una solución común cuando la red corporativa es problemática es
utilizar un \textbf{hotspot móvil} (punto de acceso creado por un teléfono), que permite
controlar directamente qué dispositivos se conectan.


\section{Sensor de temperatura y humedad DHT11}

El \textbf{DHT11} es un sensor digital de bajo coste que mide temperatura y humedad
relativa del aire. Es muy popular en proyectos educativos y de prototipado porque
su conexión es simple (solo necesita un pin de datos y alimentación) y proporciona
valores directamente en formato digital, sin necesidad de conversión analógica.

\begin{definicion}{DHT11}
El DHT11 es un sensor digital compuesto que incluye un elemento resistivo para medir
humedad y un termistor para medir temperatura. Su resolución es de 1 grado Celsius y
1 \% de humedad relativa, con un rango de 0 a 50 grados Celsius. El tiempo mínimo
entre lecturas es de aproximadamente 1 segundo.
\end{definicion}

Aunque el DHT11 no es el sensor más preciso del mercado (el DHT22 o el SHT30 son
más exactos), su simplicidad y precio lo hacen ideal para aprender los fundamentos
de la adquisición de datos y la transmisión de sensores.


\section{Práctica 3: Comunicaciones Inalámbricas}
\label{sec:practica3}

\subsection{Objetivos}

El objetivo de esta práctica es diseñar e implementar dos aplicaciones que transmitan
la temperatura ambiente del laboratorio a una red:

\begin{enumerate}
    \item Una aplicación basada en \textbf{Bluetooth BLE} que permita la comunicación
    directa entre dos placas ESP32-S3 sin cables.
    \item Una aplicación basada en \textbf{Wi-Fi} que permita acceder a la temperatura
    mediante una petición HTTP desde cualquier dispositivo conectado a la misma red.
\end{enumerate}

\subsection{Material utilizado}

\begin{itemize}
    \item \textbf{Hardware}: Ordenador, dos placas ESP32-S3-DevKitC-1 N16R8, un sensor
    de temperatura DHT11.
    \item \textbf{Software}: Arduino IDE con las librerías \texttt{DHT}, \texttt{BLEDevice}
    y \texttt{WebServer}.
    \item \textbf{Red}: Conexión a la red inalámbrica \texttt{UPV-PSK} (o un hotspot móvil
    como alternativa).
\end{itemize}

\subsection{Desarrollo de la práctica}

\subsubsection{Primera parte: lectura del sensor DHT11}

Antes de enviar datos por el aire, es necesario obtenerlos de forma fiable. El primer
programa tiene como objetivo leer la temperatura del sensor DHT11 conectado al
pin 2 de la ESP32 y mostrarla por el puerto serie.

El código es sencillo: se inicializa el sensor, se realiza una lectura cada 2 segundos
y se comprueba que el valor sea válido (el DHT11 puede devolver errores ocasionales
si no se respeta el tiempo de muestreo).

\begin{lstlisting}[caption={Lectura del sensor DHT11},label={lst:dht11-lectura}]
#include <DHT.h>

#define DHTPIN 2
#define DHTTYPE DHT11

DHT dht(DHTPIN, DHTTYPE);

void setup() {
  Serial.begin(115200);
  dht.begin();
  Serial.println("Sensor DHT11 inicializado.");
}

void loop() {
  float temperatura = dht.readTemperature();

  if (isnan(temperatura)) {
    Serial.println("Error leyendo el DHT11");
  } else {
    Serial.print("Temperatura: ");
    Serial.print(temperatura);
    Serial.println(" C");
  }
  delay(2000);
}
\end{lstlisting}

\begin{figure}[H]
\centering
\fbox{\parbox{0.8\textwidth}{\centering\vspace{1.5cm}\textit{TODO: Anadir fotografia del montaje fisico con la ESP32-S3 y el sensor DHT11 conectado.}\vspace{1.5cm}}}
\caption{Montaje fisico del sensor DHT11 con la ESP32-S3}
\label{fig:montaje-dht11}
\end{figure}

\subsubsection{Segunda parte: servidor BLE}

Una vez confirmado que el sensor funciona, se implementa un \textbf{servidor BLE}
sobre una de las ESP32-S3. Este servidor crea un servicio de temperatura con dos
características:

\begin{itemize}
    \item \textbf{Caracteristica 1}: envia la temperatura como valor numerico (float).
    \item \textbf{Caracteristica 2}: envia la temperatura como cadena de texto descriptiva.
\end{itemize}

Ambas caracteristicas tienen notificaciones activadas, lo que significa que cualquier
cliente conectado recibira los datos automaticamente sin tener que solicitarlos.

El servidor tambien gestiona la reconexion: si el cliente se desconecta, la ESP32
vuelve a anunciarse (advertising) para aceptar nuevas conexiones.

\begin{lstlisting}[caption={Servidor BLE de temperatura (extracto)},label={lst:ble-server}]
// UUIDs del servicio y caracteristicas
#define SERVICE_UUID        "8628b276-4bba-420a-ab99-dd69945408bf"
#define CHARACTERISTIC_UUID_1 "6cb1114c-4410-4757-85af-b064b3cc689a"
#define CHARACTERISTIC_UUID_2 "a1642e22-4f4e-41de-b173-b5bdb099f12c"

// Inicializar BLE
BLEDevice::init("ESP32_Temperatura");
pServer = BLEDevice::createServer();
pServer->setCallbacks(new MyServerCallbacks());

BLEService *pService = pServer->createService(SERVICE_UUID);

// Caracteristica numerica con notificaciones
pCharacteristic_1 = pService->createCharacteristic(
  CHARACTERISTIC_UUID_1,
  BLECharacteristic::PROPERTY_NOTIFY
);
pBLE2902_1 = new BLE2902();
pBLE2902_1->setNotifications(true);
pCharacteristic_1->addDescriptor(pBLE2902_1);

pService->start();
BLEDevice::startAdvertising();
\end{lstlisting}

En el bucle principal, si hay un cliente conectado, el servidor lee el DHT11 y
envia el valor por ambas caracteristicas:

\begin{lstlisting}[caption={Envio de temperatura por BLE},label={lst:ble-envio}]
if (deviceConnected) {
  float temperatura = dht.readTemperature();
  if (isnan(temperatura)) {
    Serial.println("Error leyendo el DHT11");
    delay(2000);
    return;
  }

  // Enviar como valor numerico
  pCharacteristic_1->setValue(temperatura);
  pCharacteristic_1->notify();

  // Enviar como texto descriptivo
  String mensaje = "Temperatura actual: " + String(temperatura) + " C";
  pCharacteristic_2->setValue(mensaje.c_str());
  pCharacteristic_2->notify();

  delay(2000);
}
\end{lstlisting}

\subsubsection{Tercera parte: cliente BLE}

La segunda ESP32-S3 actua como \textbf{cliente BLE}. Su mision es escanear el entorno
buscando un dispositivo que anuncie el \texttt{SERVICE\_UUID} del servidor, conectarse
a el, suscribirse a las notificaciones de una de las caracteristicas y mostrar el
dato recibido por el puerto serie.

\begin{lstlisting}[caption={Cliente BLE de temperatura (extracto)},label={lst:ble-client}]
// UUIDs (iguales que en el servidor)
#define SERVICE_UUID        "8628b276-4bba-420a-ab99-dd69945408bf"
#define CHARACTERISTIC_UUID "a1642e22-4f4e-41de-b173-b5bdb099f12c"

// Callback cuando se detecta el servidor
class MyAdvertisedDeviceCallbacks: public BLEAdvertisedDeviceCallbacks {
  void onResult(BLEAdvertisedDevice advertisedDevice) {
    if (advertisedDevice.haveServiceUUID() &&
        advertisedDevice.isAdvertisingService(BLEUUID(SERVICE_UUID))) {
      Serial.println("Servidor encontrado");
      pServerAddress = new BLEAddress(advertisedDevice.getAddress());
      doConnect = true;
      advertisedDevice.getScan()->stop();
    }
  }
};

// Callback cuando llegan notificaciones
static void notifyCallback(
  BLERemoteCharacteristic* pBLERemoteCharacteristic,
  uint8_t* pData, size_t length, bool isNotify) {
  String mensaje = "";
  for (int i = 0; i < length; i++) {
    mensaje += (char)pData[i];
  }
  Serial.print("Mensaje recibido: ");
  Serial.println(mensaje);
}
\end{lstlisting}

\begin{figure}[H]
\centering
\fbox{\parbox{0.8\textwidth}{\centering\vspace{1.5cm}\textit{TODO: Anadir captura de pantalla del monitor serie del cliente BLE mostrando la temperatura recibida.}\vspace{1.5cm}}}
\caption{Monitor serie del cliente BLE recibiendo temperatura}
\label{fig:ble-cliente-serie}
\end{figure}

\subsubsection{Cuarta parte: servidor y cliente Wi-Fi}

La alternativa a BLE es utilizar \textbf{Wi-Fi} para crear una red de sensores mas
abierta. En este caso, la ESP32-S3 que lleva el DHT11 se conecta a la red
\texttt{UPV-PSK} y actua como \textbf{servidor web HTTP} en el puerto 80.

Cada vez que un cliente visita la URL \texttt{http://IP\_DEL\_ESP32/temperatura},
el servidor lee el sensor y responde con el valor de temperatura como texto plano.

\begin{lstlisting}[caption={Servidor Wi-Fi de temperatura},label={lst:wifi-server}]
#include <WiFi.h>
#include <WebServer.h>
#include <DHT.h>

#define DHTPIN 2
#define DHTTYPE DHT11
DHT dht(DHTPIN, DHTTYPE);

const char* ssid = "UPV-PSK";
const char* password = "C0n3ct4nd0s3_m0l4!";
WebServer server(80);

void handleTemperatura() {
  float temperatura = dht.readTemperature();
  if (isnan(temperatura)) {
    server.send(500, "text/plain", "Error leyendo sensor");
    return;
  }
  server.send(200, "text/plain", String(temperatura));
}

void setup() {
  Serial.begin(115200);
  dht.begin();
  WiFi.begin(ssid, password);
  while (WiFi.status() != WL_CONNECTED) {
    delay(500);
    Serial.print(".");
  }
  Serial.println("\nConectado!");
  Serial.print("IP del servidor: ");
  Serial.println(WiFi.localIP());

  server.on("/temperatura", handleTemperatura);
  server.begin();
}

void loop() {
  server.handleClient();
}
\end{lstlisting}

Por otro lado, la segunda ESP32 (o cualquier ordenador) actua como \textbf{cliente
HTTP}. Se conecta a la misma red Wi-Fi, construye la URL a partir de la IP del
servidor, realiza una peticion \texttt{GET} cada 2 segundos y muestra la respuesta.

\begin{lstlisting}[caption={Cliente Wi-Fi de temperatura},label={lst:wifi-client}]
#include <WiFi.h>
#include <HTTPClient.h>

const char* ssid = "UPV-PSK";
const char* password = "C0n3ct4nd0s3_m0l4!";
const char* serverIP = "192.168.1.42";

void setup() {
  Serial.begin(115200);
  WiFi.begin(ssid, password);
  while (WiFi.status() != WL_CONNECTED) {
    delay(500);
    Serial.print(".");
  }
  Serial.println("\nConectado a la red WiFi");
}

void loop() {
  if (WiFi.status() == WL_CONNECTED) {
    HTTPClient http;
    String url = String("http://") + serverIP + "/temperatura";
    http.begin(url);
    int httpResponseCode = http.GET();

    if (httpResponseCode > 0) {
      String temperatura = http.getString();
      Serial.print("Temperatura recibida: ");
      Serial.println(temperatura + " C");
    } else {
      Serial.print("Error en la conexion: ");
      Serial.println(httpResponseCode);
    }
    http.end();
  }
  delay(2000);
}
\end{lstlisting}

\begin{figure}[H]
\centering
\fbox{\parbox{0.8\textwidth}{\centering\vspace{1.5cm}\textit{TODO: Anadir captura de pantalla del monitor serie del cliente Wi-Fi mostrando la temperatura recibida.}\vspace{1.5cm}}}
\caption{Monitor serie del cliente Wi-Fi recibiendo temperatura via HTTP}
\label{fig:wifi-cliente-serie}
\end{figure}

\subsection{Problemas encontrados}

Durante la practica, la conexion a la red \texttt{UPV-PSK} resulto en ocasiones
bastante complicada. Frecuentemente aparecian \textbf{timeouts} que obligaban a
reiniciar el proceso de conexion. Tras analizar la situacion, se concluyo que la red
tiene un sistema de \textbf{control de acceso} que mantiene una lista de dispositivos
admitidos. Algunos equipos quedaban en espera dentro de esta lista, lo que impedia
que se conectaran correctamente.

Como \textbf{alternativa}, se opto por utilizar el \textbf{hotspot de un telefono movil}.
Esto permitio controlar directamente que dispositivos podian conectarse, agilizando
el proceso y asegurando que la practica pudiera realizarse sin depender de las
limitaciones de la red de la universidad. Esta solucion demostro ser mas fiable y
rapida para el desarrollo de los experimentos.

\subsection{Relacion teoria-practica: BLE vs. Wi-Fi en la practica}

\begin{table}[htbp]
\centering
\small
\caption{Conexion entre conceptos teoricos y elementos de la Practica 3}
\begin{tabularx}{\textwidth}{@{}lX@{}}
\toprule
\textbf{Concepto teorico} & \textbf{Observacion/aplicacion en la practica} \\
\midrule
Bluetooth BLE & Se implemento un servidor/periférico BLE con servicio de temperatura y un cliente/central que recibe notificaciones \\
UUID (GATT) & Cada servicio y caracteristica tiene un UUID unico de 128 bits para identificarlos sin ambiguedad \\
Advertising & El servidor BLE envia paquetes de anuncio para que el cliente pueda descubrirlo \\
Notificaciones BLE & La temperatura se envia automaticamente al cliente sin que este tenga que solicitarla explicitamente \\
Wi-Fi (modo STA) & La ESP32 se conecta a la red UPV-PSK como estacion, obteniendo una IP por DHCP \\
Servidor HTTP & La ESP32 actua como servidor web en el puerto 80, respondiendo peticiones GET en la ruta /temperatura \\
Cliente HTTP & La segunda ESP32 realiza peticiones GET periodicas para obtener la temperatura del servidor \\
Red corporativa IoT & Se evidencio el problema del control de acceso en redes universitarias, resuelto con un hotspot alternativo \\
Sensor DHT11 & Proporciona el dato de temperatura real que alimenta tanto al servidor BLE como al servidor Wi-Fi \\
Consumo energetico & BLE consume menos que Wi-Fi, aunque en esta practica ambas placas estaban alimentadas por USB \\
\bottomrule
\end{tabularx}
\end{table}


\section{Resumen de la Unidad 2}

\begin{itemize}
    \item Las \textbf{comunicaciones inalambricas} eliminan la necesidad de cableado
    fisico, ofreciendo flexibilidad y movilidad, aunque introducen desafios como
    interferencias, seguridad y consumo energetico.
    \item \textbf{Bluetooth} y \textbf{Wi-Fi} comparten la banda de 2.4 GHz pero estan
    disenados para propositos distintos: BLE es ideal para sensores de bajo consumo
    a corta distancia, mientras que Wi-Fi proporciona mayor alcance y velocidad para
    acceso a redes y a internet.
    \item \textbf{BLE} organiza los datos mediante el perfil GATT, utilizando servicios
    y caracteristicas identificados por UUIDs. Las notificaciones permiten enviar
    datos de forma automatica sin polling.
    \item \textbf{Wi-Fi} en sistemas embebidos permite crear servidores HTTP que
    exponen datos de sensores a traves de peticiones web estandar, facilitando la
    integracion con aplicaciones y plataformas en la nube.
    \item La \textbf{Practica 3} demuestra la viabilidad de ambas tecnologias para
    transmitir datos de temperatura en tiempo real, comparando el enfoque punto a
    punto de BLE con el enfoque cliente-servidor web de Wi-Fi.
    \item Las \textbf{redes corporativas IoT} pueden presentar barreras de acceso
    (control de dispositivos, timeouts) que requieren soluciones alternativas como
    hotspots moviles para el desarrollo de prototipos.
\end{itemize}
